\par\Needspace{6\baselineskip}\medskip
\SourceLeaf{158}\pdfbookmark[2]{113.}{article-113}
\begin{center}113.\end{center}
To avoid being overwhelmed by the length of the formulas, we use the following abbreviations. First, we denote the quantity
\[
\operatorname{tang}\beta\sin(\alpha''-\alpha')+\operatorname{tang}\beta'\sin(\alpha-\alpha'')+\operatorname{tang}\beta''\sin(\alpha'-\alpha)
\]
by $(0.1.2)$. If the longitude and latitude corresponding to any geocentric position in that expression are replaced by the longitude and latitude corresponding to any of the Earth's three heliocentric positions, we replace the number for the former in the symbol $(0.1.2)$ with the Roman numeral for the latter. Thus, for example, the symbol $(0.1.\mathrm{I})$ represents the quantity
\[
\operatorname{tang}\beta\sin(L'-\alpha')+\operatorname{tang}\beta'\sin(\alpha-L')+\operatorname{tang}B'\sin(\alpha'-\alpha)
\]
and the symbol $(0.\mathrm{O}.2)$ represents
\[
\operatorname{tang}\beta\sin(\alpha''-L)+\operatorname{tang}B\sin(\alpha-\alpha'')+\operatorname{tang}\beta''\sin(L-\alpha).
\]
We change the symbol in the same way if, in the first expression, we replace \emph{two} geocentric pairs of longitude and latitude with any two heliocentric pairs for the Earth. If two pairs of longitude and latitude entering the same expression are merely interchanged, the corresponding numbers in the symbol must also be interchanged. The value itself then remains unchanged, merely passing from positive to negative or from negative to positive. For example, $(0.1.2)=-(0.2.1)=(1.2.0)=-(1.0.2)=(2.0.1)=-(2.1.0)$. All quantities arising in this way therefore reduce to the following 19
\[
\begin{aligned}
&(0.1.2)\\
&(0.1.\mathrm{O}),\ (0.1.\mathrm{I}),\ (0.1.\mathrm{II}),\\
&(0.\mathrm{O}.2),\ (0.\mathrm{I}.2),\ (0.\mathrm{II}.2),\\
&(\mathrm{O}.1.2),\ (\mathrm{I}.1.2),\ (\mathrm{II}.1.2)\\
&(0.\mathrm{O}.\mathrm{I}),\ (0.\mathrm{O}.\mathrm{II}),\ (0.\mathrm{I}.\mathrm{II}),\\
&(1.\mathrm{O}.\mathrm{I}),\ (1.\mathrm{O}.\mathrm{II}),\ (1.\mathrm{I}.\mathrm{II}),\\
&(2.\mathrm{O}.\mathrm{I}),\ (2.\mathrm{O}.\mathrm{II}),\ (2.\mathrm{I}.\mathrm{II}),
\end{aligned}
\]
to which a twentieth is added, $(\mathrm{O}.\mathrm{I}.\mathrm{II})$.
\par
It is readily shown that each of these expressions, when multiplied by the product of the cosines of the three latitudes entering it, equals six times the volume of the pyramid whose vertex is at the Sun and whose base is the triangle formed by the three points on the celestial sphere corresponding to the positions entering the expression, with the radius of the sphere set $=1$. Thus, whenever these three \SourceLeaf{159}positions lie on the same great circle, the value of the expression must be $=0$. This always occurs for the three heliocentric positions of the Earth whenever we disregard parallaxes and the Earth's latitudes arising from perturbations, that is, whenever we place the Earth in the plane of the ecliptic itself. With this assumption, we always have $(\mathrm{O}.\mathrm{I}.\mathrm{II})=0$. This equation is an identity if the ecliptic itself was taken as the third plane. Moreover, whenever $B$ and $B'$ vanish and $B''=0$, all these expressions except the first become much simpler. Each from the second through the tenth consists of two terms, and each from the eleventh through the nineteenth consists of a single term.
